% Mistral helped me
% \ytodo{Prolog is a natural way to "implement a language"}
Prolog is a natural choice for implementing domain-specific languages, and EMFtoCSP leverages this strength to model OCL constraints.

% \ytodo{SWI-prlog and $ECL^iPS_e$ also provide constraint programming techniques ontop of prolog techniques}
SWI-prlog and $ECL^iPS_e$ extend Prolog with constraint programming techniques, providing powerful tools for modeling and solving combinatorial problems.
% \ytodo{EMF2CSP uses $ECL^iPS_e$ to neatly describe OCL, and leverage gobal constraints such as $element$ (like us)}
EMFtoCSP specifically uses $ECL^iPS_e$ to represent OCL operations, such as sum, through global constraints like sum and element. This allows for a clean and efficient mapping of OCL constructs to constraint problems.



% \ytodo{Important note (un-like us): This method will generate a different problem for each combination of collection sizes. Whereas we add the "what size is the collection" problem to the CSP.}
However, it is important to note that EMFtoCSP generates a separate constraint problem for each combination of collection sizes.
In contrast, our approach integrates the collection size as a variable within the CSP, allowing the solver to dynamically determine the size of collections. This distinction is significant for modeling UML relations, where collection sizes can vary.

\begin{lstlisting}
    % ocl_col_sum( ?Col, ?Result ) :-
    %    Result is the sum of all elements within the collection
    %    Warning: currently this method only works properly for
    %    collections of integers
    delay ocl_col_sum(X, _) if var(X).
    ocl_col_sum(Col, Result) :- Result #= sum(Col).
\end{lstlisting}
In this listing we can see how EMFtoCSP models \inlineocl{sum} on collections of integers.
It simply employs the $sum$ global constraint\footnote{line 156 of https://github.com/SOM-Research/EMFtoCSP/blob/master/plugins/fr.inria.atlanmod.emftocsp.eclipsecs/libs/ocl_collections.ecl}.
\begin{lstlisting}
sum(Vector)
    Sum of the vector elements. 
    The vector can be a list, array or any of the vector expressions supported by eval_to_list/2.
\end{lstlisting}
In this listing we see the documentation\footnote{https://www.eclipseclp.org/doc/bips/lib/ic/index.html} for that constraint.
Telling us it takes simple arrays of variable, for which all the variables have valid values.
This indicates to us: propagation isn't being applied to the size of the list being summed upon.
Which is part of the overall problem, as UML relations can vary in size.